\section{模型的评价与总结}
\label{sec:evaluation}

\subsection{模型的优点}
三问共用统一的整图单核基准和官方事件评估器，逐步扩大接收域：问题一以Task为接收域，问题二以核心为接收域，问题三在核级计划上加入共享只读Cache的查询、完成和淘汰事件。每个图--核数--场景组合都从实际生成的候选中独立选取最终计划，主矩阵覆盖100张图、1至5核和A/B/C三种场景，共1500条官方成功记录。

问题三用同图同核配对把Cache变化拆为硬件效应、C内计划适配和综合效应，500条配对的乘积恒等式最大残差为$2.22\times10^{-16}$。五核原始配对的三段均值为$1.0146$、$1.0139$和$1.0289$；严格哈希一致的98个五核配对中，硬件改善48组、适配改善37组、综合改善64组。Cache参数敏感性中24组固定计划和37条重优化候选全部通过官方评估，重优化在四种硬件变体中均有3/6组改变计划。

\subsection{模型的不足}
不同核数的候选集合独立生成，最终选择没有强制相邻核数的工期单调性。由4核增至5核时，A、B、C分别有5、7、9张图变长；这一现象来自有限候选和每组评价预算的组合，平均速度比仍分别提高到3.4383、3.4738和3.5533。三条B/C初始计划哈希不一致记录为case\_054的5核、case\_062的4核和case\_085的5核，主矩阵与配对表保留这些记录并将其与严格同计划解释分开。

原始弱连通分量为1的16张图五核速度比均值为1.30至1.60，而超过100个分量的27张图达到4.78至4.84，说明紧密依赖链上的跨核等待和结构搬运仍限制切分收益。候选集合只覆盖有限的迁移、切分、合并、调序和共享输入操作，轻量排序与官方Makespan之间的排序一致性需要在更大候选集上继续检验。Cache零命中诊断显示，500条C记录中174条最终命中字节为零，命中候选的存在不必然带来更短尾部工期。

\subsection{算法复杂度与计算成本}
算法成本分为图索引、候选构造与轻量排序、官方完整评价三部分。令$|V|$为非COPY操作数，$|E_0|$为原始边数，$|E|$为去重后的计算依赖数，$|\mathcal T|$为张量数，$P$为生产者--消费者配对数。生产消费索引的代价为$O(|V|+|\mathcal T|+|E_0|+P)$，堆拓扑序和后继排序为$O((|V|+|E|)\log|V|)$，弱连通分量遍历为$O(|V|+|E|)$；就绪操作排序在最坏情况下为$O(|V|^2\log|V|)$。若生成$M$个候选、单候选检查与评分成本为$C_{\rm cand}$，候选阶段为$O(MC_{\rm cand}+M\log M)$。完整评价需展开Task、COPY和物理版本，其成本随展开事件数$E_{\rm ev}$与容量扫描增长。

\begin{table}[H]\centering\small
\caption{v7主矩阵最终选择的下界距离与官方评价耗时}
\label{tab:eval-cost}
\setlength{\tabcolsep}{4pt}
\begin{tabular}{lrrrrr}\toprule
场景 & 组合数 & $F/LB$中位数 & $F/LB$第90百分位 & 评价耗时中位数/s & 评价耗时最大/s\\\midrule
A & 500 & 9.556 & 35.690 & 0.433 & 14.696\\
B & 500 & 9.252 & 35.694 & 0.666 & 21.774\\
C & 500 & 9.170 & 35.493 & 0.722 & 34.064\\\bottomrule
\end{tabular}\end{table}

主矩阵最终选择中，A、B、C轻量画像时间中位数分别为0.345、0.638和0.683 s，95\%分位数分别为8.127、12.688和12.367 s；对应完整官方评价耗时中位数为0.433、0.666和0.722 s，最大值为14.696、21.774和34.064 s。官方事件评估由CPU执行，Apple MPS只用于轻量负载排序。全部主结果来自有限候选官方评价，候选生成、轻量评分和完整评价的边界均记录在实验交付包中。

\subsection{基线、消融与参数扰动}
相对整图单核基线，五核平均速度比由统一基准计算为A=3.4383、B=3.4738、C=3.5533。初始候选到最终选择的平均改善在A场景为1.7900\%至58.6697\%，B场景为4.0928\%至60.3622\%，C场景为0.8289\%至1.3230\%。消融记录覆盖48个图--核--场景组和96条候选，全部通过官方评估；候选搜索对A、B的改善主要来自结构切分与接收域调整，C的收益集中在少数Cache事件相关候选。

六图四种L2扰动的固定计划工期比均为1。重新搜索后，容量减半、容量加倍、带宽减半和带宽加倍的重优化/默认工期比均值分别为0.9900、0.9869、0.9892和0.9885；四种变体均有3/6组改变计划，37条重优化候选全部通过官方评估。容量加倍的平均命中率由固定计划的8.11\%提高到重优化的13.02\%，容量减半时为2.10\%和3.44\%；带宽两种扰动的重优化命中率为7.28\%和7.35\%。计划改变集中在\Case{016}、\Case{062}和\Case{067}，说明L2参数对切图调度的影响通过容量压力、读请求时序和尾部关键路径共同体现。

\subsection{模型推广}
模型适用于依赖图、操作成本和片上容量已知，且能够对完整计划执行离散事件评价的多核调度任务。结构搬运公式刻画接收域的复制代价，物理副本生存区间对应容量换出，FIFO查询与完成事件决定共享读服务。实际部署时，可按图规模、候选展开事件数和$F/LB$距离分配评价预算，并对大图设置候选上限；这一调整只改变搜索资源分配，不改变官方事件评价规则。

\subsection{全文结论}
本文建立了从Task接收域切分、核心级分配到Cache感知事件调度的三层模型。v7主矩阵在100张图、1至5核和A/B/C三场景上完成1500条官方评价，五核平均速度比分别为3.4383、3.4738和3.5533。B/C配对表明，Cache硬件效应和C内计划适配分别提供约1.46\%和1.39\%的五核平均工期比增益；L2敏感性实验进一步显示，重新搜索后容量与带宽扰动均能在少数容量压力图上改变计划并缩短工期。结果给出了切图、核内顺序和Cache参数之间的可计算联系，也保留了有限候选、三条计划配对差异和独立预算试验等实验边界。
